\documentclass[10pt,a4paper]{amsart}
\usepackage[margin=19mm]{geometry}
\usepackage{amsmath,amssymb,mathtools}
\usepackage[scr=rsfs,cal=euler]{mathalfa}
\usepackage{xfrac}
\usepackage[unicode,hidelinks,bookmarks=false]{hyperref}
\newcommand{\ie}{i.e.\ }
\newcommand{\cf}{cf.\ }
\newcommand{\resp}{respectively\ }
\newcommand{\Subsection}{\subsection}
\providecommand{\nobreakdash}{\nobreak\hbox{-}}
\setlength{\emergencystretch}{2em}
\multlinegap0pt
\usepackage{bm}
\usepackage{bbm}
\usepackage{dsfont}
\usepackage{xspace}
\usepackage{cancel}
\usepackage{mathtools}
\usepackage{amsthm}
\makeatletter
\@ifundefined{lemma}{\newtheorem{lemma}{Lemma}}{}
\@ifundefined{proposition}{\newtheorem{proposition}[lemma]{Proposition}}{}
\makeatother
\makeatletter
\expandafter\ifx\csname proof\endcsname\relax
\renewenvironment{proof}[1][\proofname]{\par
  \pushQED{\qed}%
  \normalfont \topsep6\p@\@plus6\p@\relax
  \trivlist
  \item[\hskip\labelsep
        \itshape
    #1\ \ref{#1}.]\ignorespaces
}{%
  \popQED\endtrivlist\@endpefalse
}
\fi
\makeatother
\title[Birth control and turnpike property of Lotka-McKendrick mode]{Birth control and turnpike property of Lotka-McKendrick models with diffusion}
\author{Marius Bargo, Yacouba Simpore}
\date{Source: arXiv:2409.11247v2 (CC BY 4.0)}
\begin{document}
\maketitle

\noindent\emph{Source and licence.} Birth control and turnpike property of Lotka-McKendrick models with diffusion. Version arXiv:2409.11247v2 (CC BY 4.0), distributed under the Creative Commons Attribution 4.0 International licence, as recorded in the arXiv licence metadata for this exact version and shown on the abstract page https://arxiv.org/abs/2409.11247v2. Full rights notice: licenses/arxiv-2409.11247-CC-BY-4.0.txt.\par\smallskip

\noindent\textit{Editorial scope.} Counted unit: the Proposition whose source label is \texttt{pr2.3} in the article (source0.tex lines 2469--2475, source label \texttt{pr2.3}), delivered together with the article's own complete original proof of it (source0.tex lines 2476--2641, one \texttt{proof} environment of 166 lines). No mathematics is written, repaired or re-derived: statement and proof are the article's own text line for line. Required context retained uncounted: eq1.5 (lines 720--729); th1.3 (lines 731--740); eq2.6 (lines 975--982). Every definition, display and label of those lines is retained unchanged. Omitted from this package and not redistributed: 1--719, 730--730, 741--974, 983--2468, 2642--2991. Nothing inside a retained excerpt is omitted or altered: every retained body is delivered exactly as the article's lines are written, so no omission receipt of any kind accompanies this package. Citations: the retained excerpts carry no citation command at all, so no bibliography entry of the article is transcribed and no reference is redistributed. Reference adaptation: none was needed and none was made; every reference inside the delivered unit resolves inside it, so no unresolved reference or citation is produced. Printed slips kept literal: no printed word, symbol, hypothesis or number of the counted statement or of its proof was altered. Layout and class: the article's own class is \texttt{amsart} with packages including lmodern, inputenc, amsmath, graphicx, amssymb, esint, xcolor, tikz, xxcolor, floatrow, color, amsthm, epsfig, babel; the wrapper is the lane's pinned \texttt{amsart} editorial wrapper at 10pt on A4, which restates the theorem-like environments the retained text needs and adds the article's own macro definitions that the retained text needs (none), with extra wrapper packages (bm, bbm, dsfont, xspace, cancel, mathtools). The delivered numbering is the unit's own. Assets: no retained span contains a figure environment, a graphics-inclusion command, a PDF-page inclusion, a table or a TikZ picture; no image file is copied into this package, the package declares no asset dependency and no image file is redistributed with it. Licence: version arXiv:2409.11247v2 of this article is distributed under the Creative Commons Attribution 4.0 International licence, as recorded in the arXiv licence metadata for this exact version and shown on the abstract page https://arxiv.org/abs/2409.11247v2. A full rights notice is delivered at licenses/arxiv-2409.11247-CC-BY-4.0.txt. Characters: every retained excerpt is pure ASCII, so the delivered engine, XeLaTeX with its default Latin Modern text font, reports no missing character.
\begin{equation}\label{eq1.5}
\left\lbrace 
\begin{array}{ll}
\partial_{t}y(x,a,t)+\partial_{a}y(x,a,t)-\triangle y(x,a,t)+\mu(a) y(x,a,t)=\mathds{1}_{\lbrace\omega\times[0,a_0]\rbrace}(x,a)v(x,a,t) &\text{in }\quad Q,\\
 \\\partial_{\nu}y(x,a,t)=0 &\text{on }\quad \Sigma,\\
 \\y(x,0,t)=\displaystyle\int_{0}^{A}\beta(x,a)y(x,a,t)da &\text{in }\quad Q_{T},\\
 \\y(x,a,0)=y_{0} &\text{in }\quad Q_{A},
\end{array}
\right.
\end{equation} 

\begin{proposition}\label{th1.3}
Consider the assumptions (H1-H2). Furthermore, suppose that the fertility rate is such that
\begin{equation*}
\beta(x,a)=0\qquad\qquad \forall\quad a\in [0, a_b]
\end{equation*}
for $a_b \in [0, A]$ and $a_b>0$. Thus, for every $T>A-a_0$ and  $y_0\in H,$ there exists a control $v\in L^2(\omega\times(0,a_0)\times(0,T) )$ such that $y$ solution to \eqref{eq1.5} satisfies
\begin{equation}\label{e3.2}
y(x,a,T)=0\qquad x\in\Omega,\,a\in [0,A],\quad\text{a.e.}.
\end{equation}
\end{proposition}

\begin{equation}\label{eq2.6}
y(x,a,t)=\left\lbrace
\begin{array}{ll}
\frac{\pi(a)}{\pi(a-t)}e^{t\triangle}y_{0}(x,a-t)+u(x,t-a)\displaystyle\int_{a-t}^{a}\frac{\pi(a)}{\pi(s)}e^{(a-s)\triangle}\mathds{1}_{\lbrace\omega\times[0,a_0]\rbrace}(x,s)ds\quad t\leq a,\quad \forall x\in \Omega\\
\\\pi(a)e^{a\triangle}b(x,t-a)+u(x,t-a)\displaystyle\int_{0}^{a}\frac{\pi(a)}{\pi(s)}e^{(a-s)\triangle}\mathds{1}_{\lbrace\omega\times[0,a_0]\rbrace}(x,s)ds\quad t> a, \quad \forall x\in \Omega.
\end{array}
\right.
\end{equation}

\begin{proposition}\label{pr2.3}  
Let us define the admissible set as follows:  
\begin{equation*}
\mathcal{R}=\left\lbrace \bar{y}\in H \mid \dfrac{ e^{-a\triangle}\bar{y}(x,a)}{\pi(a)}\in H ,\quad \dfrac{\dfrac{e^{-a\triangle}\bar{y}(x,a)}{\pi(a) }-\displaystyle\int_{0}^{A-a}\beta(x,z)\dfrac{\pi(z)}{\pi(z+a)}e^{-a\triangle}\bar{y}(x,z+a)dz}{\displaystyle\int_{0}^{a}\frac{e^{-z\triangle}\mathds{1}_{\lbrace\omega\times[0,a_0]\rbrace}(x,z)}{\pi(z) }dz}\in H \right\rbrace.
\end{equation*}  
Under the assumptions of Proposition \ref{th1.3}, we have that \(\mathcal{R}\subset \operatorname{Ran}(\Phi_T)\) for every \(T > A - a_0\).
\end{proposition}

\begin{proof}[of Proposition \ref{pr2.3}]
Before determining the admissible space, we first investigate the existence of a control. We assume that \( a_0 \leq a_b \) and consider \( T \in \;] A-a_0, A] \). Our objective is to find a control function \( u(x,t-a) \in L^2 (\Omega\times [-A,T]) \) satisfying  
\begin{equation}\label{eq2.8}
(\Phi_T v)(a) =\bar{y}(x,a) \qquad \forall \;a\in [0,A],\quad x\in \Omega.
\end{equation}  

{\bf \underline{Case 1 : \( -A < t-a < T-A \)}}
We first set \( u(x,t-a) = 0 \) and, using equation \eqref{eq2.6}, we obtain  
\begin{equation}\label{eq2.9}
(\Phi_T v)(a) = 0 \qquad \forall \;t-a<T-A.
\end{equation}  

{\bf \underline{Case 2: \( T-A \leq t-a \leq 0 \)}}  
By employing equation \eqref{eq2.6}, we derive  
\begin{equation}
(\Phi_T v)(a) = u(x,T-a) \int_{a-T}^{a} \frac{\pi(a)}{\pi(s)} e^{(a-s)\triangle} \mathds{1}_{\lbrace\omega\times[0,a_0]\rbrace}(x,s) ds, \qquad \forall a\in [T,A],\quad \forall x\in \Omega.
\end{equation}  

To ensure that equation \eqref{eq2.8} holds, we must solve  
\begin{align*}
u(x,T-a) \int_{a-T}^{a} \frac{\pi(a)}{\pi(s)} e^{(a-s)\triangle} \mathds{1}_{\lbrace\omega\times[0,a_0]\rbrace}(x,s) ds &= \bar{y}(x,a), \qquad \forall a\in [T,A], \\
u(x,T-a) &= \dfrac{e^{-a\triangle} \bar{y}(x,a) / \pi(a)}{\displaystyle\int_{a-T}^{a} \frac{e^{-s\triangle} \mathds{1}_{\lbrace\omega\times[0,a_0]\rbrace}(x,s)}{\pi(s)} ds}.
\end{align*}  

By setting \( T-a = s \), we rewrite the control as  
\begin{equation}\label{eq2.11}
u(x,s) = \dfrac{e^{(s-T)\triangle} \bar{y}(x,T-s) / \pi(T-s)}{\displaystyle\int_{-s}^{T-s} \frac{e^{-z\triangle} \mathds{1}_{\lbrace\omega\times[0,a_0]\rbrace}(x,z)}{\pi(z)} dz}, \qquad \forall s\in [T-A,0],\quad \forall x\in \Omega.
\end{equation}  

Since we have \( -a_0 < T-A \leq s \), it follows that \( -s < a_0 \), which ensures that  
\begin{equation*}
\int_{-s}^{T-s} \frac{e^{-z\triangle} \mathds{1}_{\lbrace\omega\times[0,a_0]\rbrace}(x,z)}{\pi(z)} dz > 0.
\end{equation*}  
 \underline{\bf{Case 3 : $0\leq t-a\leq  a_b$}} 
%To establish the existence of a control, we employ an inductive method, as the renewal equation retains a trace at \( t-a \), thereby depending on the initial control. Indeed, u
Using \eqref{eq2.6} and leveraging \eqref{eq2.11}, we obtain  
\begin{align}\label{eq2.12}
(\Phi_t v)(a) = \pi(a) \dfrac{\displaystyle\int_{a-t}^{a} \frac{e^{-z\triangle} \mathds{1}_{\lbrace\omega\times[0,a_0]\rbrace}(x,z)}{\pi(z)} dz}{\displaystyle\int_{a-t}^{T+(a-t)} \frac{e^{-z\triangle} \mathds{1}_{\lbrace\omega\times[0,a_0]\rbrace}(x,z)}{\pi(z)} dz} \dfrac{e^{(t-T)\triangle} \bar{y}(x,T+(a-t))}{\pi(T+(a-t))}.
\end{align}  

%### **Determination of the Control over \( [0, a_b] \)**  
From equation \eqref{eq2.6}, we have  
\begin{equation}\label{eq2.13}
(\Phi_T v)(a) = \pi(a)e^{a\triangle}b(x,T-a) + u(x,T-a) \int_{0}^{a} \frac{\pi(a)}{\pi(s)} e^{(a-s)\triangle} \mathds{1}_{\lbrace\omega\times[0,a_0]\rbrace}(x,s) ds, \quad a\in [ T-a_b, T],
\end{equation}  
where  
\begin{equation*}
b(x,t) = \int_{a_b}^{A} \beta(x,a)(\Phi_t u)(a) da, \qquad t\in [0,a_b].
\end{equation*}  
Decomposing \( b(x,t) \), we write  
\begin{align*}
b(x,t) &= \int_{a_b}^{t+A-T} \beta(x,a)(\Phi_t u)(a) da + \int_{t+A-T}^{A} \beta(x,a)(\Phi_t u)(a) da, \\
b(x,t) &= \int_{a_b}^{t+A-T} \beta(x,a)(\Phi_t u)(a) da.
\end{align*}  
For \( a \leq t+A-T \), from \eqref{eq2.12}, we obtain  
\begin{equation}
b(x,t) = \int_{a_b}^{t+A-T} \beta(x,a)  \pi(a) \dfrac{\displaystyle\int_{a-t}^{a} \frac{e^{-z\triangle} \mathds{1}_{\lbrace\omega\times[0,a_0]\rbrace}(x,z)}{\pi(z)} dz}{\displaystyle\int_{a-t}^{T+(a-t)} \frac{e^{-z\triangle} \mathds{1}_{\lbrace\omega\times[0,a_0]\rbrace}(x,z)}{\pi(z)} dz} \dfrac{e^{(t-T)\triangle} \bar{y}(x,T+(a-t))}{\pi(T+(a-t))} da.
\end{equation}  
Moreover,  
\begin{equation*}
\int_{a-t}^{T+(a-t)} \frac{e^{-z\triangle} \mathds{1}_{\lbrace\omega\times[0,a_0]\rbrace}(x,z)}{\pi(z)} dz = \int_{a-t}^{a} \frac{e^{-z\triangle} \mathds{1}_{\lbrace\omega\times[0,a_0]\rbrace}(x,z)}{\pi(z)} dz.
\end{equation*}  
Thus,  
\begin{equation*}
b(x,t) = \int_{a_b}^{t+A-T} \beta(x,a) \pi(a) \dfrac{e^{(t-T)\triangle} \bar{y}(x,T+(a-t))}{\pi(T+(a-t))} da, \qquad t\in [0,a_b], \quad \forall x\in \Omega,
\end{equation*}  
which implies  
\begin{equation*}
b(x,T-a) = \int_{a_b}^{A-a} \beta(x,z)  \pi(z) \dfrac{e^{-a\triangle} \bar{y}(x,z+a)}{\pi(z+a)} dz, \qquad a\in[T-a_b, T].
\end{equation*}  

Substituting this result into \eqref{eq2.13}, we obtain  
\begin{equation}
(\Phi_T v)(a) = \pi(a) \int_{a_b}^{A-a} \beta(x,z)  \pi(z) \dfrac{\bar{y}(x,z+a)}{\pi(z+a)} dz + u(x,T-a) \pi(a) \int_{0}^{a} \frac{e^{(a-s)\triangle} \mathds{1}_{\lbrace\omega\times[0,a_0]\rbrace}(x,s)}{\pi(s)} ds, \quad a\in [T-a_b, T].
\end{equation}  

To satisfy \eqref{eq2.8}, we impose  
\begin{align*}
 \pi(a) \int_{a_b}^{A-a} \beta(x,z)  \pi(z) \dfrac{\bar{y}(x,z+a)}{\pi(z+a)} dz + u(x,T-a) \pi(a) \int_{0}^{a} \frac{e^{(a-s)\triangle} \mathds{1}_{\lbrace\omega\times[0,a_0]\rbrace}(x,s)}{\pi(s)} ds = \bar{y}(x,a).
\end{align*}  
This leads to the control function  
\begin{align}
    u(x,T-a) = \dfrac{e^{-a\triangle} \dfrac{\bar{y}(x,a)}{\pi(a)} - \displaystyle\int_{a_b}^{A-a} \beta(x,z)  \pi(z) \dfrac{e^{-a\triangle} \bar{y}(x,z+a)}{\pi(z+a)} dz}{\displaystyle\int_{0}^{a} \frac{e^{-s\triangle} \mathds{1}_{\lbrace\omega\times[0,a_0]\rbrace}(x,s)}{\pi(s)} ds}, \quad a\in [T-a_b,T],
\end{align}  
where  
\begin{equation*}
\int_{0}^{a} \frac{e^{-s\triangle} \mathds{1}_{\lbrace\omega\times[0,a_0]\rbrace}(x,s)}{\pi(s)} ds > 0.
\end{equation*}  

By setting \( s = T-a \), we further refine  
\begin{align*}
u(x,s) = \dfrac{\dfrac{e^{(s-T)\triangle} \bar{y}(x,T-s)}{\pi(T-s)} - \displaystyle\int_{a_b}^{A-T+s} \beta(x,z)  \pi(z) \dfrac{e^{(s-T)\triangle} \bar{y}(x,T-(s-z))}{\pi(T-(s-z))} dz}{\displaystyle\int_{0}^{T-s} \frac{e^{-z\triangle} \mathds{1}_{\lbrace\omega\times[0,a_0]\rbrace}(x,z)}{\pi(z)} dz}, \quad s\in[0,a_b].
\end{align*}  

Defining  
\begin{equation}
\chi(x,s) = \int_{a_b}^{A-T+s} \beta(x,z)  \pi(z) \dfrac{e^{(s-T)\triangle} \bar{y}(x,T-(s-z))}{\pi(T-(s-z))} dz, \quad s\in[0,a_b],
\end{equation}  
we obtain the final expression  
\begin{align}\label{eq2.17}
    u(x,s) = \dfrac{e^{(s-T)\triangle} \bar{y}(x,T-s) - \pi(T-s) \chi(x,s)}{\pi(T-s) \displaystyle\int_{0}^{T-s} \frac{e^{-z\triangle} \mathds{1}_{\lbrace\omega\times[0,a_0]\rbrace}(x,z)}{\pi(z)} dz}, \quad s\in[0,a_b].
\end{align}  
\underline{\bf{Case 4 : $a_b\leq t-a\leq T$}}  
In this section, we will prove that  
\begin{align}
b(x,t) =\chi (x,t), \qquad t\in [ja_b, (j+1)a_b],
\end{align}  
and that \eqref{eq2.8} holds for all \( a \in [T-(j+1)a_b, T-ja_b] \) with \( j \geq 0 \).  

Indeed, from \eqref{eq2.17}, we have  
\begin{align*}
(\Phi_t v)(a) &= \pi(a)e^{a\triangle}b(x,t-a) + u(x,t-a) \int_{0}^{a} \frac{\pi(a)}{\pi(s)} e^{(a-s)\triangle} \mathds{1}_{\lbrace\omega\times[0,a_0]\rbrace}(x,s) ds, \\
(\Phi_t v)(a) &= \pi(a)e^{a\triangle}b(x,t-a) + \frac{e^{((t-a)-T)\triangle} \bar{y}(x,T-(t-a)) - \pi(T-(t-a)) \chi(x,t-a)}{\pi(T-(t-a)) \int_{0}^{T-(t-a)} \frac{e^{-z\triangle} \mathds{1}_{\lbrace\omega\times[0,a_0]\rbrace}(x,z)}{\pi(z)} dz} \\
&\quad \times \int_{0}^{a} \frac{\pi(a)}{\pi(z)} e^{(a-z)\triangle} \mathds{1}_{\lbrace\omega\times[0,a_0]\rbrace}(x,z) dz, \\
(\Phi_t v)(a) &= \pi(a)e^{a\triangle}b(x,t-a) \left[1 - \frac{\int_{0}^{a} \frac{e^{-z\triangle} m(x,z)}{\pi(z)} dz}{\int_{0}^{T-(t-a)} \frac{e^{-z\triangle} \mathds{1}_{\lbrace\omega\times[0,a_0]\rbrace}(x,z)}{\pi(z)} dz} \right] \\
&\quad + \pi(a) \frac{\int_{0}^{a} \frac{e^{-z\triangle} m(x,z)}{\pi(z)} dz}{\int_{0}^{T-(t-a)} \frac{e^{-z\triangle} \mathds{1}_{\lbrace\omega\times[0,a_0]\rbrace}(x,z)}{\pi(z)} dz} \times \frac{e^{(t-T)\triangle} \bar{y}(x,T-(t-a))}{\pi(T-(t-a))}.
\end{align*}  

To compute \( b(x,t) \) over the interval \([(j+1)a_b,(j+2)a_b]\), we use the fact that  
\begin{align*}
b(x,t) = \int_{a_b}^{t+A-T} \beta(x,a) (\Phi_t u)(a) da.
\end{align*}  
Thus, we decompose it as follows:  
\begin{align}
b(x,t) = \int_{a_b}^{t} \beta(x,a) (\Phi_t u)(a) da + \int_{t}^{t+A-T} \beta(x,a) (\Phi_t u)(a) da.
\end{align}  
After computation, we obtain  
\begin{align}
b(x,t) = \int_{a_b}^{A-T+t} \beta(x,z) \pi(z) \frac{e^{(t-T)\triangle} \bar{y}(x,T-(t-z))}{\pi(T-(t-z))} dz, \qquad t\in[(j+1)a_b,(j+2)a_b].
\end{align}  

In summary, the control function is given by  
\begin{equation}
u(x,s) =
\begin{cases}
0, & s \in [-A,T-A], \quad \forall x \in \Omega, \\
\frac{\frac{e^{(s-T)\triangle} \bar{y}(x,T-s)}{\pi(T-s)}}{\int_{-s}^{T-s} \frac{e^{-z\triangle} \mathds{1}_{\lbrace\omega\times[0,a_0]\rbrace}(x,z)}{\pi(z)} dz}, & s \in [T-A,0], \quad \forall x \in \Omega, \\
\frac{e^{(s-T)\triangle} \bar{y}(x,T-s) - \pi(T-s) \chi(x,s)}{\pi(T-s) \int_{0}^{T-s} \frac{e^{-z\triangle} \mathds{1}_{\lbrace\omega\times[0,a_0]\rbrace}(x,z)}{\pi(z)} dz}, & s \in [0,T], \quad \forall x \in \Omega,
\end{cases}
\end{equation}  
satisfying \eqref{eq2.8}, where  
\begin{equation}
\chi(x,t) = \int_{a_b}^{A-T+t} \beta(x,z) \pi(z) \frac{e^{(t-T)\triangle} \bar{y}(x,T-(t-z))}{\pi(T-(t-z))} dz, \qquad t \in [0,T], \quad \forall x \in \Omega.
\end{equation}  

Consequently, we have  
\begin{equation}
v(x,a,t) =
\begin{cases}
0, & t-a<T-A, \quad \forall x \in \Omega, \\
\frac{\frac{e^{((t-a)-T)\triangle} \bar{y}(x,T-(t-a))}{\pi(T-(t-a))}}{\int_{a-t}^{T-(t-a)} \frac{e^{-z\triangle} \mathds{1}_{\lbrace\omega\times[0,a_0]\rbrace}(x,z)}{\pi(z)} dz}, & T-A \leq t-a \leq 0, \quad \forall x \in \Omega, \\
\frac{e^{((t-a)-T)\triangle} \bar{y}(x,T-(t-a)) - \pi(T-(t-a)) \chi(x,t-a)}{\pi(T-(t-a)) \int_{0}^{T-(t-a)} \frac{e^{-z\triangle} \mathds{1}_{\lbrace\omega\times[0,a_0]\rbrace}(x,z)}{\pi(z)} dz}, & 0 < t-a \leq T, \quad \forall x \in \Omega.
\end{cases}
\end{equation}  

Finally, the operator \(\Phi_t v\) satisfies  
\begin{equation}
(\Phi_t v)(a) =
\begin{cases}
0, & t-a<T-A, \quad \forall x \in \Omega, \\
\pi(a) \frac{\int_{a-t}^{a} \frac{e^{-z\triangle} \mathds{1}_{\lbrace\omega\times[0,a_0]\rbrace}(x,z)}{\pi(z)} dz}{\int_{a-t}^{T+(a-t)} \frac{e^{-z\triangle} \mathds{1}_{\lbrace\omega\times[0,a_0]\rbrace}(x,z)}{\pi(z)} dz} \frac{e^{(t-T)\triangle} \bar{y}(x,T+(a-t))}{\pi(T+(a-t))}, & T-A \leq t-a \leq 0, \quad \forall x \in \Omega, \\
\pi(a)e^{a\triangle}b(x,t-a) \left[1 - \frac{\int_{0}^{a} \frac{e^{-z\triangle} \mathds{1}_{\lbrace\omega\times[0,a_0]\rbrace}(x,z)}{\pi(z)} dz}{\int_{0}^{T-(t-a)} \frac{e^{-z\triangle} \mathds{1}_{\lbrace\omega\times[0,a_0]\rbrace}(x,z)}{\pi(z)} dz} \right] \\
+ \pi(a) \frac{\int_{0}^{a} \frac{e^{-z\triangle} \mathds{1}_{\lbrace\omega\times[0,a_0]\rbrace}(x,z)}{\pi(z)} dz}{\int_{0}^{T-(t-a)} \frac{e^{-z\triangle} \mathds{1}_{\lbrace\omega\times[0,a_0]\rbrace}(x,z)}{\pi(z)} dz} \frac{e^{(t-T)\triangle} \bar{y}(x,T-(t-a))}{\pi(T-(t-a))}, & 0 < t-a \leq T, \quad \forall x \in \Omega.
\end{cases}
\end{equation}  
\end{proof}

\end{document}
